\section{Evaluation Method}
\label{sec:evaluation}

We executed all 760 runs of a frozen plan on one tested source snapshot
(version 2.1.29), with 20 predeclared seeds per condition. Every key
completed on its first attempt and passed the formal evidence audit.
Development pilots are excluded from these performance estimates.

\subsection{Comparison scope}

Table~\ref{tab:models} summarizes the implemented models. All use the
same custom packet encoding, radio model, application traffic, and
measurement deadline. The encoding does not reproduce product encryption,
authentication, or complete wire formats. The Meshtastic-like model is
inspired by managed flooding~\cite{meshtasticdocs}. The Meshtastic-like and
MeshCore-like models omit application retransmission, although the
reviewed official implementations or documentation describe retry and
flood fallback~\cite{meshtasticfirmware,meshcore}. The Meshtastic source
reference pins release v2.7.26.54e0d8d; the MeshCore FAQ pins companion-v1.18.0.
These are documented antecedents, not the implementations executed by
the simulator. Their limitations are part
of the interpretation of each result, not resolved by parameter matching.

\begin{table}[t]
\caption{Implemented behavioral models. These are not product firmware runs.}
\label{tab:models}
\centering
\small
\setlength{\tabcolsep}{3pt}
\begin{tabular}{@{}>{\raggedright\arraybackslash}p{0.19\columnwidth}>{\raggedright\arraybackslash}p{0.36\columnwidth}>{\raggedright\arraybackslash}p{0.38\columnwidth}@{}}
\toprule
Model & Initial route acquisition & Subsequent source behavior \\
\midrule
MeshEcho-CPR & Useful DATA flood; ACK confirms path & Cached route, one repeat, one admitted recovery flood \\
Legacy MeshEcho & RREQ/RREP with candidate collection & Cached route; limited discovery/fallback policy \\
Meshtastic-like & Useful DATA flood per message & Duplicate suppression; no application retransmission \\
MeshCore-like & RREQ/RREP; first reply gives route & Cached route; no application retransmission \\
\bottomrule
\end{tabular}
\end{table}

Two regimes are reported separately. The implemented-default regime
uses each model's resolved defaults under shared PHY and workload settings.
The matched regime aligns only parameters with comparable meaning:
MeshCore-like route TTL becomes 600~s, and Meshtastic-like and MeshCore-like
flood base delay/jitter become 0.45/0.65~s. The corresponding defaults are
300~s for the MeshCore-like TTL, 0.75/0.75~s for Meshtastic-like delays,
and 0.5/0.7~s for MeshCore-like delays. CPR and legacy MeshEcho already
use 600~s and 0.45/0.65~s. SNR-dependent delay functions, route-learning
semantics, discovery collection, and retry policies remain distinct.
Equal base delays do not imply equal total forwarding delay.

Each clean CPR control retains the reference's resolved parameters except
its named switch. ACK-only cancellation, repeat, and recovery controls are
defined in Section~\ref{sec:cpr}. Any additional differing effective
parameter makes a comparison a mixed-policy comparison; it cannot enter
the clean ablation panel.

\subsection{Common environment and scenario validity}

\begin{table}[t]
\caption{Frozen common environment and finite scenario axes.}
\label{tab:settings}
\centering\small
\setlength{\tabcolsep}{3pt}
\begin{tabular}{@{}>{\raggedright\arraybackslash}p{0.32\columnwidth}>{\raggedright\arraybackslash}p{0.63\columnwidth}@{}}
\toprule
Parameter & Value \\
\midrule
Topology & 20 nodes, $5\times4$ grid, 1000-m spacing; position jitter fraction 0.025 \\
Propagation & Path-loss exponent 2.75; extra loss 20 dB for neighbors, 80 dB otherwise \\
Random channel & Shadowing / temporal fading SD 2 / 2 dB; 60-s fading interval; 0.01-s channel bins \\
Radio / application & 915 MHz, SF7, 125 kHz, coding rate 4/5; 17 dBm; 6-dB capture threshold; 32-byte payload \\
Endpoints / hop limit & $5\leftrightarrow9$, $10\leftrightarrow14$; $H=6$ \\
Duration / deadline & 240-s traffic generation; 30-s deadline and drain in main cases \\
S1 & No link fault; 8 offered flows/min, whole network \\
S2 & S1 plus 80-dB attenuation on edges 7--8 and 12--13 in $[80,160)$ s \\
S3 & No link fault; 40 offered flows/min, whole network \\
Deadline response & S2 matched, 10 / 30 / 60-s deadline and corresponding drain \\
\bottomrule
\end{tabular}
\end{table}

Table~\ref{tab:settings} specifies the common environment. Traffic uses
Poisson arrivals over four directed endpoint pairs. The grid's fault-free
shortest paths have four hops; six-hop structural alternatives avoid the
two S2 fault edges. These are synthetic stress conditions, not a calibrated
deployment. The two fault-free loads reuse S1 and S3; they are two observed
points, not a saturation-threshold sweep.

The scenario generator supplies shared topology, application arrivals,
endpoint pairs, and exogenous channel/fault fields. Their content hashes
match within each paired case and seed. Channel outcomes use a
link/time-bin field, avoiding dependence on how many prior random draws
another protocol consumed. This is a controlled coupling of the model's
environment; protocols can still transmit at different times, encounter
different interference, and obtain different outcomes. Common random
numbers do not guarantee variance reduction for every system
\cite{glasserman1992crn}.

Scenario interpretation uses observed initial-route choices, actual
fault-edge transmissions, repeat/recovery starts, and contention/state
diagnostics. Fault exposure depends on each transmission's start time;
mere interval overlap is insufficient. Structural reachability is an
offline diagnostic, not a guarantee of decoding or ACK return. All flows
remain in the scheduled denominator.

\subsection{Frozen population and statistical quantities}

A run identity binds cohort, configuration, source snapshot, case, arm,
and seed. The plan contains 480 main runs (three settings, two regimes,
four models, 20 seeds), 120 component-control runs (three controls in S2/S3
matched), and 160 extra deadline-response runs (10/60 s, four models,
20 seeds). The 30-s response reuses S2 matched. The fixed corrected
controller was retained after development without selecting a new
candidate on formal outcomes. Seeds 2027100--2027119 were registered as
a separate validation cohort before formal execution; all are retained.
The sample size was fixed, with no outcome selection or precision stopping.
Default/matched duplicates are reported separately and do not increase
the independent seed count to 40.

Source/configuration hashes, completion records, shared-input identities,
and physical-event joins passed audit before aggregation. Independent
reviews checked the full raw population, statistics, and exposed paths.
Source, plan, tables, event hashes, and correction history are retained
in the reproducibility package.

The reproducibility record includes a post-start administrative erratum
clarifying stale development-stage headings after 43 runs had completed;
the original freeze, scientific settings, seed cohort, and 760-key plan
remained unchanged. Execution continued with the remaining registered keys,
without a partial performance summary or parameter selection at that pause.

Each seed, rather than each packet, is the analysis unit. For a metric
$X$, each arm's estimate is $n^{-1}\sum_s X_s$. For comparator $c$,
the paired confirmation effect and airtime ratio are
\begin{equation}
\begin{aligned}
\overline{\delta P_A}&=\frac{1}{n}\sum_s
 (P_{A,s}^{\mathrm{CPR}}-P_{A,s}^{c}),\\
G_A&=\exp\left\{\frac{1}{n}\sum_s
\log\left(\frac{A_s^{\mathrm{CPR}}}{A_s^c}\right)\right\}.
\end{aligned}
\label{eq:paired}
\end{equation}
The second quantity is a geometric mean of seed-level ratios, not a ratio
of pooled means. The analysis uses 10,000 paired seed resamples and
percentile 95\% intervals. These intervals are descriptive and unadjusted
for multiple comparisons; they do not establish confirmatory significance,
noninferiority, or universal dominance. The present $n=20$ supports finite,
descriptive replication of the chosen scenarios rather than precise
characterization of all operating conditions.

Generation lasts 240~s in every deadline case; the drain equals its
10-, 30-, or 60-s deadline. This axis therefore changes the accounting
horizon together with the application constraint. The plan includes no
cross-deadline inferential contrast. Connected points describe the finite
response and must not be read as a continuous fitted relationship.

An undefined seed-level quantity remains undefined in its all-seed
summary; zero timely deliveries do not yield zero cost per delivery.
Success-only latency accompanies the unfinished fraction. The packet/stage
airtime partition reconstructs the total. ACK checks, successful
cancellations, and capacity evictions remain distinct; emitted events
cannot reconstruct exact peak ledger occupancy.
